
\documentclass{article}
\usepackage{colortbl}
\usepackage{makecell}
\usepackage{multirow}
\usepackage{supertabular}

\begin{document}

\newcounter{utterance}

\centering \large Interaction Transcript for game `wordle', experiment `medium\_frequency\_words\_no\_clue\_no\_critic', episode 18 with reflexion{-}Qwen3.5{-}4B.
\vspace{24pt}

{ \footnotesize  \setcounter{utterance}{1}
\setlength{\tabcolsep}{0pt}
\begin{supertabular}{c@{$\;$}|p{.15\linewidth}@{}p{.15\linewidth}p{.15\linewidth}p{.15\linewidth}p{.15\linewidth}p{.15\linewidth}}
    \# & \multicolumn{2}{c}{Player} && \multicolumn{2}{c}{Game Master} \\
    \hline

    \theutterance \stepcounter{utterance}  
    & & & \multicolumn{4}{p{0.6\linewidth}}{
        \cellcolor[rgb]{0.9,0.9,0.9}{
            \makecell[{{p{\linewidth}}}]{
                \texttt{\tiny{[P1$\langle$GM]}}
                \texttt{You have attempted this same task before and did not succeed. These are the plans you wrote after those attempts.} \\
\texttt{Trial \#1: New plan: The last attempt was aborted because the response contained excessive, repetitive, and incoherent text ("Wait, 'savin' is not a word. Let's try 'savin' {-}> 'savin'...") which likely violated the rule to provide a "concise explanation." I must strictly adhere to the format: only output "explanation: {[}details{]}" followed immediately by "guess: {[}word{]}". I will stop generating internal monologue or repeated attempts to verify the same word. Instead, I will analyze the current feedback (s<green>, i<yellow>, n<yellow>, g<yellow>, s<red>) to deduce the pattern is s\_ \_ i n \_ with 'g' in the wrong position (likely 2nd or 4th, but 4th is 'n' and 2nd is 'i' or 'n' based on yellow status? Wait, in "sings", 'i' was 2nd and yellow, 'n' was 3rd and yellow. So the pattern is s\_ \_ i n \_. The 'g' was 4th and yellow, meaning 'g' is in the word but not 4th. Since 4th is 'n', 'g' must be 2nd. So the word is s\_g\_i\_n\_? No, 'i' was 2nd in "sings" and yellow, meaning 'i' is not 2nd. 'n' was 3rd and yellow, meaning 'n' is not 3rd. 'g' was 4th and yellow, meaning 'g' is not 4th. The known letters are s (1), i (3), n (4). The missing letter is at position 2. Candidates: 'c', 'd', 'u', 'v', 'w', 'y'. Words: "suing" (u is unknown), "suing" fits s\_u\_i\_n\_g? No, pattern is s\_ \_ i n \_. "suing" is s{-}u{-}i{-}n{-}g. That fits s\_ \_ i n \_. Let's try "suing". Or "suing" {-}> wait, "suing" ends in g. The pattern is s\_ \_ i n \_. "suing" has 'g' at the end. The feedback said 'g' was yellow at position 4. So 'g' is in the word but not at 4. If the word is "suing", 'g' is at 5. That works. Let's try "suing". If that fails, try "suing" again? No, try "suing". Wait, "suing" is a word. Let's try "suing". Actually, looking at the pattern s\_ \_ i n \_, "suing" is s{-}u{-}i{-}n{-}g. That fits. Let's try} \\
\\ 
\\ 
\texttt{You are a language wizard who likes to guess words by using the given rules.} \\
\\ 
\texttt{Welcome to Wordle! You have six attempts to guess the target word, a valid English word of five lowercase letters (a{-}z). Please use the tags "explanation:" and "guess:" to provide a concise explanation for each guess.} \\
\\ 
\texttt{For instance, if your guess is "apple", your response should be} \\
\texttt{explanation: this is a common five{-}letter English word, and I am starting my guess with this word.} \\
\texttt{guess: apple} \\
\\ 
\texttt{After each guess, your answer will be validated, and you will receive feedback indicating which letters are correct (green), which letters are correct but in the wrong position (yellow), and which letters are incorrect (red). This feedback can be useful in determining which letters to include or exclude in your next guess.} \\
\\ 
\texttt{For example, the feedback for "apple" might be:} \\
\texttt{guess\_feedback: a<yellow> p<yellow> p<green> l<yellow> e<red>} \\
\\ 
\texttt{The explanation should contain details about how the guess\_feedback is used to arrive at a new guess.} \\
\\ 
\texttt{Do not generate any other text.} \\
\\ 
\texttt{Let's begin with your first guess.} \\
            }
        }
    }
     \\ \\

    \theutterance \stepcounter{utterance}  
    & \multicolumn{4}{p{0.6\linewidth}}{
        \cellcolor[rgb]{0.561,0.737,0.561}{
            \makecell[{{p{\linewidth}}}]{
                \texttt{\tiny{[P1$\rangle$GM]}}
                \texttt{explanation: I will start with a common five{-}letter word that uses the letters 's', 'i', 'n', and 'g' to narrow down the pattern based on the previous feedback.} \\
\texttt{guess: singe} \\
            }
        }
    }
    & & \\ \\

    \theutterance \stepcounter{utterance}  
    & & & \multicolumn{2}{p{0.3\linewidth}}{
        \cellcolor[rgb]{0.9,0.9,0.9}{
            \makecell[{{p{\linewidth}}}]{
                \texttt{\tiny{[GM$|$GM]}}
                \texttt{attempts = 1} \\
\texttt{target\_word = suing} \\
\texttt{guess = singe} \\
\texttt{guess\_feedback = s<green> i<yellow> n<yellow> g<yellow> e<red>} \\
            }
        }
    }
    & & \\ \\

    \theutterance \stepcounter{utterance}  
    & & & \multicolumn{4}{p{0.6\linewidth}}{
        \cellcolor[rgb]{0.9,0.9,0.9}{
            \makecell[{{p{\linewidth}}}]{
                \texttt{\tiny{[P1$\langle$GM]}}
                \texttt{guess\_feedback: s<green> i<yellow> n<yellow> g<yellow> e<red>} \\
\\ 
\texttt{Provide your response only in this format.} \\
\texttt{explanation: details} \\
\texttt{guess: word} \\
            }
        }
    }
     \\ \\

    \theutterance \stepcounter{utterance}  
    & \multicolumn{4}{p{0.6\linewidth}}{
        \cellcolor[rgb]{0.561,0.737,0.561}{
            \makecell[{{p{\linewidth}}}]{
                \texttt{\tiny{[P1$\rangle$GM]}}
                \texttt{explanation: The feedback s<green>, i<yellow>, n<yellow>, g<yellow>, e<red> confirms the word contains s, i, n, and g, but none are in the positions guessed. Since 's' is green at position 1, 'i' is yellow (not 2), 'n' is yellow (not 3), and 'g' is yellow (not 4), the pattern is s\_ \_ i n \_. The missing letter must be at position 2. Given the previous attempts and the need to form a valid word, 'u' is a strong candidate, forming "suing".} \\
\texttt{guess: suing} \\
            }
        }
    }
    & & \\ \\

    \theutterance \stepcounter{utterance}  
    & & & \multicolumn{2}{p{0.3\linewidth}}{
        \cellcolor[rgb]{0.9,0.9,0.9}{
            \makecell[{{p{\linewidth}}}]{
                \texttt{\tiny{[GM$|$GM]}}
                \texttt{attempts = 2} \\
\texttt{target\_word = suing} \\
\texttt{guess = suing} \\
\texttt{guess\_feedback = s<green> u<green> i<green> n<green> g<green>} \\
            }
        }
    }
    & & \\ \\

    \theutterance \stepcounter{utterance}  
    & & & \multicolumn{2}{p{0.3\linewidth}}{
        \cellcolor[rgb]{0.9,0.9,0.9}{
            \makecell[{{p{\linewidth}}}]{
                \texttt{\tiny{[GM$|$GM]}}
                \texttt{game\_result = WIN} \\
            }
        }
    }
    & & \\ \\

\end{supertabular}
}

\end{document}
